\pdfminorversion=7
\documentclass[opre,dblanonrev]{informs4}
\RequirePackage{tgtermes}
\RequirePackage{newtxtext}
\RequirePackage{newtxmath}
\RequirePackage{bm}
\OneAndAHalfSpacedXI
\usepackage{mathtools}
\usepackage{booktabs}
\usepackage{multirow}
\usepackage{microtype}
\usepackage{natbib}
\bibpunct[, ]{(}{)}{,}{a}{}{,}
\def\newblock{\ }
\def\BIBand{and}
\usepackage{xr-hyper}
\usepackage[hidelinks]{hyperref}
\externaldocument[][nocite]{meta_prior_main}[meta_prior_main.pdf]
\EquationsNumberedThrough
\TheoremsNumberedThrough
\ECRepeatTheorems
\setlength{\emergencystretch}{2em}

\ECRUNAUTHOR{Authors' names not included for peer review}
\ECRUNTITLE{Source-Learned Structural Priors for Chance Constraints}
\hypersetup{pdftitle={Online Supplement: Source-Learned Structural Priors for Chance-Constrained Simulation Optimization},pdfauthor={}}

\begin{document}
\ECSwitch
\raggedbottom
\begin{center}
{\Large\bfseries Online Supplement}\par
\vspace{0.6em}
{\large Source-Learned Structural Priors for Chance-Constrained
Simulation Optimization}\par
\vspace{0.4em}
Authors' names are not included for peer review
\end{center}

\section{Frozen Protocol and Accounting}

\begin{table}[h]
\centering
\caption{The primary source-to-target information contract.}
\label{tab:mp_protocol}
\begin{tabular}{ll}
\toprule
Source tasks & 8\\
Public candidates per source task & 64\\
Source repetitions per candidate & 12\\
Target initial points & 8\\
Target search points & 24\\
Independent verification repetitions & 928\\
Mean/risk heads & $f$, $g$, $\eta=\log\operatorname{Var}(Y^g)$\\
Dictionary size & 68 centered/scaled functions\\
Original confirmation interval & not run as a block; seed 95434 overlaps development\\
\bottomrule
\end{tabular}
\end{table}

Source calls are 8$\times$64$\times$12$=6{,}144$.  The paired
source-structure control uses two model arms but the same keyed target stream;
its additional local execution made 4,608 simulator calls for the two arms and
eight tasks.  These are paired development calls, not independent confirmation
observations.  Oracle moments are read only after recommendation for scoring and
are never supplied to fitting or acquisition.

\subsection{Public coordinates and the finite dictionary}

Each candidate is an ordered inventory policy
$p=(p_1,p_2,p_3,p_4)\in[0,1]^4$.  With capacity $C=20$ and initial stock
$s_0=8$, the observed stock bin is
$b_t=\min\{\lfloor 4s_t/C\rfloor,3\}$, the order is
$o_t=\max\{Cp_{b_t+1}-s_t,0\}$, and the next stock is
$s_{t+1}=\max\{s_t+o_t-d_t,0\}$.  For the declared twelve-period demand schedule,
the local trajectory weights are
\[
w_{t1}=o_t/(C T),\qquad w_{t2}=s_{t+1}/(C T),\qquad
w_{t3}=\max\{d_t-s_t-o_t,0\}/\sum_t d_t.
\]
Let $A_j=\|w_{\cdot j}\|_2$ for $j=1,2,3$ and
$(N_1,N_2)=(\sum_t w_{t1},\,\sum_t w_{t3})$ for the
two shared volume exposures used by the simulator.  The public coordinate is
\[
\psi(p)=(A_1,A_2,A_3,N_1,N_2,h,\nu,s_T/C),
\]
where $h=\sum_t w_{t2}$, $\nu=T^{-1}\sum_t\mathbb{1}\{o_t>0\}$, and
$s_T$ is terminal stock.  Thus the eight coordinates are observable from the
trajectory and do not contain prices, outcomes, or a hidden optimum.

The implementation uses a 68-function dictionary.  After pooled source
standardization, it contains the eight coordinates $\psi_j$, their eight squares,
the sixteen $\sin(\pi\psi_j)$ and $\cos(\pi\psi_j)$ terms, the 28 pairwise
products $\psi_i\psi_j$ $(i<j)$, and the eight high-frequency controls
$\sin(4\pi\tanh(\psi_j/2))$.  The source learner scores these fixed functions by
low-frequency smoothness, cross-task relevance, sign consistency, and a pilot
cross-validation loss, then greedily retains four functions per response head.
The full-dictionary arm retains all 68 functions; the target-only arm uses the
same 68 functions with a zero-mean coefficient prior and no source labels.
No trajectory encoder or learned input coordinate is used.

\subsection{Source prior, target update, and verifier}

For each mean response head $t\in\{f,g\}$, source labels are fitted by
weighted ridge regression on the selected standardized functions.  The source
task coefficient vectors supply a head-specific mean and covariance; the target
coefficient mean starts at zero in the target-only control.  Paid target labels
update the Gaussian linear model with the observed diagonal noise.  The
physical-risk head additionally updates a finite factor-family particle mixture
by the exact Gaussian log-sample-variance likelihood.  The risk head is the total log variance
$\eta=\log\operatorname{Var}(Y^g)$; its simulator law is
\[
\operatorname{Var}(Y^g)=\rho_0+A^\top\Lambda A+N^\top B N+N^\top\omega.
\]
The observation at a policy uses 12 independent episodes.  Its objective and
constraint means have noise variances $\sigma_f^2/12$ and $\widehat v_g/12$;
the log-sample-variance label uses the digamma bias correction and noise
$\psi_1((12-1)/2)$.

The acquisition rule uses a three-dimensional Gaussian approximation for
$(f,g,\eta)$ and a 3-node Hermite rule per head, hence 27 joint fantasy nodes.
For a candidate $x$, the reported deployment utility is
\[
U(x)=\Pr\{g(x)+z_{.95}\exp(\eta(x)/2)\le 0\}\,
\mathbb{E}\left[\frac{5-f(x)}{10}\right],
\]
with objective quality clipped to the declared $[-5,5]$ simulator range.  The
928-repetition verifier is independent of paid observations and returns zero
realized utility when its chance certificate fails.  These definitions make the
target-only baseline and the source-selected arm identical in candidate bank,
budget, noise streams, acquisition, and verification.

\section{Source-Structure Control}

Table~\ref{tab:mp_taskwise} gives the complete paired final-reward comparison
for the eight development tasks.  Both arms use observable inventory exposure,
the same pooled factor-family risk prior, the same initial points, the same
Gaussian KG rule, and the same verifier.

\begin{table}[h]
\centering
\caption{Taskwise source-selected structure minus same-source full dictionary.}
\label{tab:mp_taskwise}
\begin{tabular}{rrrrrr}
\toprule
Law seed & Initial reward & Structure & Full dictionary & Difference & Both certified\\
\midrule
95425 & 0.875956 & 0.881686 & 0.881686 & +0.000000 & yes\\
95426 & 0.880274 & 0.880274 & 0.878995 & +0.001279 & yes\\
95427 & 0.785842 & 0.859288 & 0.859288 & +0.000000 & yes\\
95428 & 0.855284 & 0.871915 & 0.873167 & -0.001253 & yes\\
95430 & 0.824105 & 0.851198 & 0.851198 & +0.000000 & yes\\
95431 & 0.795239 & 0.795239 & 0.751917 & +0.043322 & yes\\
95432 & 0.848587 & 0.848587 & 0.848340 & +0.000247 & yes\\
95434 & 0.898958 & 0.901946 & 0.901946 & +0.000000 & yes\\
\midrule
Mean & & 0.861267 & 0.855817 & +0.005449 & 8/8\\
\bottomrule
\end{tabular}

\end{table}

The positive-difference count is 3/8, the negative count is 1/8, and four
tasks are tied.  The development contract therefore records an attribution
signal but fails the preregistered stable-deployment gate.

\section{Risk and Representation Diagnostics}

The observable coordinate reduces eight-point holdout objective-mean RMSE from
0.594 to 0.270 and constraint-mean RMSE from 0.204 to 0.141.  Constraint-rank
Spearman correlation increases from 0.119 to 0.624, while the initial true
recommendation is feasible in 7/8 tasks for both coordinate choices.

The factor-family risk prior reduces initial target $\eta$ RMSE from 0.363 to
0.301 on all eight development tasks.  Its coefficient RMSE is 0.255 versus
0.251 for pooled risk, and its closed-loop mean final reward is 0.8599 versus
0.8605.  These numbers separate total-risk prediction from factor recovery and
 deployment value.

\section{Same-Contract Target-Only Baseline}

The target-only full-dictionary arm uses the same 68-function basis, target
initial design, paid-observation budget, keyed noise streams, pooled risk arm,
acquisition rule, and independent verifier.  Its only difference is that it
starts without source labels.  The complete comparison is:

\begin{table}[h]
\centering
\caption{Same-contract target-only baseline.  Parentheses give certified tasks.}
\label{tab:mp_target_only}
\begin{tabular}{lrrrrr}
\toprule
Block & Source-selected & Full dictionary & Target-only & RBF-KG & Source--RBF\\
 & mean (certified) & mean (certified) & mean (certified) & mean (certified) & mean difference\\
\midrule
Development & 0.86127 (8/8) & 0.85582 (8/8) & 0.76199 (7/8) & 0.86239 (8/8) & -0.00112 \\
Independent v2 & 0.44274 (5/8) & 0.45160 (4/8) & 0.45148 (4/8) & 0.45912 (4/8) & -0.01638 \\
\bottomrule
\end{tabular}

\end{table}

On the development block, source-selected structure, same-source full
dictionary, and target-only full dictionary obtain mean final rewards 0.86127,
0.85582, and 0.76199, respectively.  On the independent v2 block the three
means are 0.44274, 0.45160, and 0.45148.  The target-only arm therefore provides
the required feasible nonempty comparator, while the independent block still
does not support a stable source-structure deployment advantage.

For a generic constrained-BO control, we additionally froze an isotropic RBF
kernel (lengthscale 0.25, outputscale 1.0) and fit all three heads only to
target observations.  This target-only RBF-KG arm obtains means 0.86239 and
0.45912 on development and v2, respectively.  It is therefore a stronger
generic control for decision ranking than the source-selected arm on both
blocks; the source-structure attribution remains a comparison with the
same-source full dictionary, not a claim of generic BO superiority.

\section{Paired Uncertainty Summaries}

All taskwise differences reported in the manuscript are accompanied by a
seed-level paired percentile bootstrap with 100,000 resamples.  These intervals
are descriptive because each block has eight law seeds; they are not a test of
universal superiority.

\begin{table}[h]
\centering
\caption{Paired bootstrap intervals for reported differences.}
\label{tab:mp_uncertainty}
\scriptsize
\begin{tabular}{p{0.56\textwidth}r r l}
\toprule
Metric & $n$ & Mean difference & Percentile bootstrap 95\% interval\\
\midrule
observable minus policy objective RMSE & 8 & -0.32418 & [-0.48330, -0.16914] \\
observable minus policy constraint RMSE & 8 & -0.06390 & [-0.12552, -0.00427] \\
observable minus policy log-risk RMSE & 8 & +0.00000 & [+0.00000, +0.00000] \\
observable minus policy constraint rank correlation & 8 & +0.50491 & [+0.19041, +0.77954] \\
learned minus pooled log-risk RMSE & 8 & -0.06202 & [-0.08830, -0.04120] \\
learned minus pooled factor coefficient RMSE & 8 & +0.00439 & [-0.03358, +0.03654] \\
development: source\_selected\_structure\_minus\_same\_source\_full\_dictionary & 8 & +0.00545 & [-0.00031, +0.01641] \\
development: source\_selected\_structure\_minus\_target\_only\_full\_dictionary & 8 & +0.09928 & [-0.00044, +0.29825] \\
development: same\_source\_full\_dictionary\_minus\_target\_only\_full\_dictionary & 8 & +0.09383 & [-0.00048, +0.28197] \\
development: source\_selected\_structure\_minus\_target\_only\_rbf\_kg & 8 & -0.00112 & [-0.00476, +0.00129] \\
development: same\_source\_full\_dictionary\_minus\_target\_only\_rbf\_kg & 8 & -0.00657 & [-0.02106, +0.00126] \\
development: target\_only\_full\_dictionary\_minus\_target\_only\_rbf\_kg & 8 & -0.10040 & [-0.30293, +0.00136] \\
independent\_v2: source\_selected\_structure\_minus\_same\_source\_full\_dictionary & 8 & -0.00886 & [-0.02879, +0.00184] \\
independent\_v2: source\_selected\_structure\_minus\_target\_only\_full\_dictionary & 8 & -0.00873 & [-0.02879, +0.00197] \\
independent\_v2: same\_source\_full\_dictionary\_minus\_target\_only\_full\_dictionary & 8 & +0.00012 & [-0.00002, +0.00039] \\
independent\_v2: source\_selected\_structure\_minus\_target\_only\_rbf\_kg & 8 & -0.01638 & [-0.04255, +0.00738] \\
independent\_v2: same\_source\_full\_dictionary\_minus\_target\_only\_rbf\_kg & 8 & -0.00752 & [-0.02714, +0.00939] \\
independent\_v2: target\_only\_full\_dictionary\_minus\_target\_only\_rbf\_kg & 8 & -0.00765 & [-0.02726, +0.00933] \\
\bottomrule
\end{tabular}

\end{table}

\section{Claim Audit}

\begin{table}[h]
\centering
\caption{Permitted interpretation of each main claim.}
\label{tab:mp_claim_audit}
\begin{tabular}{p{0.27\textwidth}p{0.27\textwidth}p{0.34\textwidth}}
\toprule
Claim & Evidence status & Manuscript wording\\
\midrule
Source function structure improves early prediction & supported & Claim within the declared representation and dictionary.\\
Target observations correct source mismatch & supported & State that component weights can be downweighted.\\
Risk structure is learned & partial & Report total-risk calibration only.\\
Risk factors are separately recovered & unsupported & Keep outside contributions.\\
Stable independent deployment advantage & unsupported & Report the failed 3/8 gate and negative controls.\\
General high-dimensional free-variable optimization & unsupported & Describe ordered profiles and nominal dimension only.\\
Zero false certificates are universal safety & unsupported & Restrict the statement to the declared verifier law.\\
\bottomrule
\end{tabular}
\end{table}

\section{Reproducibility Pointers}

The source archive and development artifact are in
\texttt{paper\_artifacts/}\allowbreak\texttt{meta\_prior\_feasible\_deployment\_v1/}.
The paired control is regenerated by
\texttt{performance/}\allowbreak\texttt{run\_meta\_prior\_feasible\_source\_arm\_v1.py}
from \texttt{performance/}\allowbreak\texttt{meta\_prior\_feasible\_source\_arm\_v1.json}.
The same-contract target-only controls are regenerated from
\texttt{performance/}\allowbreak\texttt{meta\_prior\_feasible\_baseline\_development\_v1.json}
and \texttt{performance/}\allowbreak\texttt{meta\_prior\_feasible\_baseline\_v1.json}.
The baseline table and paired bootstrap interval table are generated by
\texttt{performance/render\_meta\_prior\_baseline.py} and
\texttt{performance/summarize\_meta\_prior\_uncertainty.py}; neither script
starts a simulator.
The original confirmation interval was not run as a block.  Seed 95434 overlaps
the development artifact and is excluded from any clean-confirmation claim.
The separate v2 block is regenerated from its own
frozen contract and retains all eight new seeds; it is reported as an
independent bounded check, not as a replacement for the original block.  The
legacy atlas manuscript and its supplement retain their own artifact identity
and are not substituted for this source-to-target protocol.

\section{Independent Fixed-Seed Block}

The v2 block uses law seeds 95800--95807, with the same public candidate bank,
paired keyed target streams, and independent verification procedure.  The
oracle preflight was recorded before target observations and did not remove a
seed.  The generated table below is read directly from the frozen result
artifact.

\begin{table}[t]
\centering
\caption{Independent fixed-seed confirmation block; all eight tasks retained.}
\begin{tabular}{lrrrr}
\toprule
Arm & Tasks & Mean final reward & Certified & False certificates\\
\midrule
Source-selected structure & 8 & 0.44274 & 5/8 & 0\\
Same-source full dictionary & 8 & 0.45160 & 4/8 & 0\\
\bottomrule
\end{tabular}
\end{table}
\par\smallskip
All eight fixed seeds were retained after the oracle-only preflight; paired
final-reward differences were positive on 2, negative on 2, and tied on 4 tasks.


\section{Retained Frontend and Native-Transfer Evidence}

The historical component study uses the three fixed families FactorShock,
InventorySupplyChain, and QueueResourceControl, with 20 algorithmic seeds
per family.  Its initial policies were frozen before target observations.
The following common-frontend backend table retains the eight supplied-frontend
transfer methods.  The native table retains each method's own initialization.
Both use the historical v69 terminal protocol; the search and source budgets
are separate from terminal verification calls.  Their totals are not pooled
with the current physical-Inventory mean/risk study.

\begin{table}[htbp]
\centering
\small
\caption{Historical supplied-frontend backend comparison. Every transfer method
receives the frozen frontend, 384 source calls and 13 target search calls.}
\label{tab:mp_legacy_backend}
\resizebox{\textwidth}{!}{\begin{tabular}{lrrrrr}
\toprule
Method & Feasible & Certified & False & Med. regret & Mean verify \\
\midrule
Atlas + canonical SAASBO & 60/60 & 60/60 & 0 & 0.0038 & 173.9 \\
Safe F-PACOH & 56/60 & 56/60 & 0 & 0.0082 & 215.5 \\
RGPE & 59/60 & 59/60 & 0 & 0.0082 & 154.9 \\
Stacked GP & 60/60 & 60/60 & 0 & 0.0082 & 90.9 \\
Transfer MTGP & 57/60 & 57/60 & 0 & 0.0082 & 204.3 \\
FSBO & 59/60 & 59/60 & 0 & 0.0082 & 181.6 \\
HyperBO & 60/60 & 60/60 & 0 & 0.0082 & 211.7 \\
MetaBO & 60/60 & 60/60 & 0 & 0.0082 & 158.9 \\
MALIBO & 60/60 & 60/60 & 0 & 0.0082 & 162.9 \\
\bottomrule
\end{tabular}
}
\end{table}

\begin{table}[htbp]
\centering
\small
\caption{Historical native transfer comparison. Each method retains its
own initialization and adaptation with the same source and search budgets.}
\label{tab:mp_legacy_native}
\resizebox{0.92\textwidth}{!}{\begin{tabular}{lrrrrr}
\toprule
Native method & FactorShock & Inventory & Queue & Total & False \\
\midrule
Safe F-PACOH & 0/20 & 0/20 & 0/20 & 0/60 & 0 \\
RGPE & 14/20 & 16/20 & 0/20 & 30/60 & 0 \\
Stacked GP & 0/20 & 13/20 & 0/20 & 13/60 & 0 \\
MTGP & 1/20 & 0/20 & 0/20 & 1/60 & 0 \\
FSBO & 6/20 & 2/20 & 0/20 & 8/60 & 0 \\
HyperBO & 0/20 & 0/20 & 0/20 & 0/60 & 0 \\
MetaBO & 19/20 & 9/20 & 0/20 & 28/60 & 0 \\
MALIBO & 12/20 & 2/20 & 0/20 & 14/60 & 0 \\
\bottomrule
\end{tabular}
}
\end{table}

The compact source summaries are
\path{paper_artifacts/paper_result_audit_summary_v1.csv} and
\path{paper_artifacts/or_review/native_transfer.json}.
\path{performance/render_legacy_frontend_evidence.py} generates the
aggregate frontend table and its evidence mapping without starting a
simulator.  The detailed frontend factorial is retained at
\path{tables/frontend_backend.tex}.  The method versions are recorded in
\path{performance/manifests/crossdim_frozen_proposal_backend_matrix_v1.json}
and \path{performance/manifests/v9_equal_terminal_verifier_crossdim_gate_v1.json}.

The development integration check in the main text is run by
\path{performance/run_legacy_joint_bridge.py}; its fixed contract is
\path{performance/manifests/legacy_joint_bridge_v1.json}.
The compact extracted source observations and seed-80 designs, per-arm paid
trajectories, solver status and terminal outcomes are retained under
\path{paper_artifacts/legacy_joint_bridge_v1}.  The original 120-iteration run
is retained separately from the completed 512-iteration source fit; the latter
changes the numerical iteration allowance, with unchanged likelihood and
optimization tolerances.  All nine source heads converge.  Each complete run
uses 48 actual search episodes, including shared initial observations, and
2,096 terminal episodes across the recorded stages.  Thus the two runs use
96 search and 4,192 terminal episodes in total, with zero new source calls.
Post-decision diagnosis is reproduced by
\path{performance/analyze_legacy_joint_bridge.py}, which makes no new calls.
These are development runs on three fixed laws with one seed per law.

\end{document}
